\documentclass[12pt]{article}

% Packages
\usepackage[margin=1in]{geometry}
\usepackage{amsmath}
\usepackage{booktabs}
\usepackage{hyperref}
\usepackage{natbib}
\usepackage{setspace}
\usepackage{parskip}
\usepackage{graphicx}
\usepackage{caption}
\usepackage{subcaption}
\usepackage{array}
\usepackage{multirow}

% Line spacing
\doublespacing

% Title block
\title{Hip Fracture Risk Prediction in Diabetic Patients:\\
A Comparison of BEHRT-LoRA, XGBoost, and Neural Network Models}
\author{Anna Gerasimenko}
\date{\today}

\begin{document}

\maketitle

\begin{abstract}
TODO: Abstract text here.
\end{abstract}

\newpage



%% ============================================================
%%  METHODS SECTION TEMPLATE
%%  Hip Fracture Risk Prediction in Diabetic Patients:
%%  BEHRT-LoRA vs. XGBoost vs. Neural Network
%% ============================================================

\section{Methods}

%% ----------------------------------------------------------
\subsection{Study Design and Overview}
%% ----------------------------------------------------------

In this retrospective cohort study, we used the IBM MarketScan
Commercial Claims Database to develop and internally validate three machine
learning models to predict hip fracture risk in adult patients with
diabetes.
The prediction task is framed as four-class classification, stratifying
patients by the interval between their first recorded diabetes diagnosis and
their first recorded hip fracture.
%% TODO: Add more details about the models we used and their architectures.

Trained models were subsequently applied without re-training to an
independent external validation cohort drawn from the All of Us Research
Program. 
This allowed us to assess generalizability across diverse populations.


%% ----------------------------------------------------------
\subsection{Data Sources}
%% ----------------------------------------------------------

\subsubsection{IBM MarketScan Commercial Claims Database}

%% TODO: Insert years of data, enrollment figures, and any
%%       relevant database version information here.
MarketScan contains de-identified administrative claims from
employer-sponsored health plans across the United States.
Data elements used include inpatient and outpatient diagnosis codes
(ICD-9-CM and ICD-10-CM), enrollment records, and patient
demographic information.

\subsubsection{All of Us Research Program}

%% TODO: Insert CDR version, enrollment figures, and any IRB
%%       or data-access agreement language required by the journal.
The All of Us Research Program is a longitudinal prospective cohort
that enrolls participants from across the United States, with the goal
of reflecting the diversity of the U.S.\ population.
EHR data are provided through Controlled Tier access to the
Researcher Workbench (CDR~v8; controlled tier).
For this study, only participants with at least one EHR condition
record were included.

%% ----------------------------------------------------------
\subsection{Study Cohort and Inclusion Criteria}
%% ----------------------------------------------------------

Patients were identified from the IBM MarketScan database by applying three
sequential inclusion criteria.

\paragraph{Enrollment duration.}
Patients were required to be enrolled for at least one full year of continuous insurance coverage.
This criterion yielded $N=$138{,}879{,}229 continuously enrolled patients.

\paragraph{Adult age at enrollment.}
Only patients aged 20 years or older at the start of their first enrollment
interval were retained. This criterion yielded $N=$179{,}415{,}164 patients.

\paragraph{Diabetes diagnosis.}
Patients were required to have at least one claim mapped to a
diabetes-specific PhecodeX code: \texttt{EM\_202} (type~2 diabetes mellitus),
\texttt{EM\_202.1}, \texttt{EM\_202.2}, or \texttt{BI\_181} (type~1 diabetes
mellitus).
This criterion yielded $N=$18{,}110{,}999 patients.

With all three criteria applied, $N=$12{,}174{,}785 patients remained.
These patients formed the final study cohort.

\paragraph{Outcome definition and risk-group assignment.}

Some patients had multiple diabetes diagnoses. The ealiest diabetes diagnosis date was defined as the index diabetes date for each patient. 
PhecodeX codes \texttt{MS\_745.11} (fracture of femur) and \texttt{MS\_745.9} (pathological fracture) were used to identify patients with a hip fracture and the earliest hip fracture date was assigned as the outcome index date for each patient.

We then split the target cohort into four mutually exclusive risk groups based on the number of days between the diabetes index date and the outcome index date.

\begin{table}[h]
\centering
\caption{Risk-group definitions and cohort counts (MarketScan).}
\label{tab:groups}
\begin{tabular}{clp{6cm}r}
\toprule
Group & Label & Definition & $N$ \\
\midrule
0 & Control         & No hip fracture, fracture on or before diabetes index date,
                      or fracture $>$10 years after diabetes & 12{,}055{,}663 \\
1 & $<$1 year       & Hip fracture 1--365 days after diabetes index date & 37{,}646 \\
2 & 1--5 years      & Hip fracture 366--1{,}825 days after diabetes index date & 58{,}478 \\
3 & 5--10 years     & Hip fracture 1{,}826--3{,}650 days after diabetes index date & 22{,}998 \\
\bottomrule
\end{tabular}
\end{table}


%% ----------------------------------------------------------
\subsection{Feature Representation}
%% ----------------------------------------------------------

\subsubsection{BEHRT: EHR Token Sequences}

\paragraph{Code representation.}
We started with raw ICD-9-CM and ICD-10-CM diagnosis codes from the MarketScan claims and mapped them to PhecodeX phenotypes. 
This mapping collapses semantically related ICD codes into a single clinically phenotype label. 

\paragraph{Sequence construction.}
For each patient, all diagnosis events were sorted chronologically and restricted to those occurring before the hip fracture outcome index date.
Hip fracture phenotypes (\texttt{MS\_745.11} and \texttt{MS\_745.9}) were excluded entirely to prevent outcome leakage.
Events occurring on the same calendar day were grouped together; a separator marker (\texttt{SEP}) was placed between consecutive visit days, and a final \texttt{SEP} was appended at the end of the record.
Age at each event was computed as an integer year from the patient's birth year and the event date.
Sex was recorded as a single patient-level value and attached to every event in the record, since it does not vary over time.

\paragraph{Model input format.}
BEHRT requires each feature to be expressed as a numeric vector rather than a raw label or number.
Each phenotype code, integer age, and sex value is therefore looked up in a learned table that maps it to a fixed-length vector of real numbers --- an \textit{embedding} --- so that the model can represent and compare them numerically.
The phenotype embedding table was built from all unique PhecodeX strings observed in the MarketScan corpus plus five special tokens (\texttt{PAD}, \texttt{CLS}, \texttt{SEP}, \texttt{MASK}, \texttt{UNK}), yielding 3{,}309 entries in total.
Integer ages were mapped to whole-year bins spanning 0 to 119, giving an age table of 122 entries (including \texttt{PAD} and \texttt{UNK}).
Sex used three entries: \texttt{PAD}, male, and female.
Two additional structural embeddings are also added at each position: a \textit{segment embedding} that alternates between two values every time a visit boundary (\texttt{SEP}) is crossed, telling the model which visit a given event belongs to; and a \textit{position embedding} that increments at the visit level rather than the token level, encoding the order of visits across the patient's history.
All five embeddings --- phenotype, age, sex, segment, and position --- are summed into a single vector at each position, and the model receives the full patient history as an ordered sequence of these per-event vectors.

The transformer can only process 512 positions at once, a fixed architectural limit.
When a patient's sequence exceeds that, the oldest events are dropped and only the most recent 511 are kept, on the assumption that recent clinical history is most predictive.
One position is always reserved for a special \texttt{CLS} token prepended to the front of every sequence, hence 511 rather than 512.
The \texttt{CLS} token has no clinical meaning; it acts as a collection point that the transformer uses to aggregate information from the entire sequence.
By the time the sequence exits the transformer, the vector at the \texttt{CLS} position has attended to every other position and effectively summarises the whole patient history --- that single vector is then passed to the classification head to predict the risk class.


\paragraph{All of Us Research Validation Cohort}


\subsubsection{XGBoost and Neural Network: One-Hot Encoded Feature Vectors}

For the gradient-boosted decision tree (XGBoost) and feedforward neural network models, each patient's clinical history was represented as a binary feature vector. No chronological information was preserved.

The feature space consisted of $N=$3{,}304 unique PhecodeX phenotypes observed across the MarketScan training cohort. 
Each element in the vector served as a binary indicator, set to 1 if the corresponding phenotype was present in the patient's clinical record and 0 otherwise. 
To rigorously prevent outcome leakage, a temporal cutoff was applied: for patients in the three positive risk groups (hip fracture within 10 years), only phenotype codes recorded strictly before the outcome index date were included in the feature vector. 
For patients in the control group, no temporal restriction was applied. 
Furthermore, the target hip fracture phenotypes (\texttt{MS\_745.11} and \texttt{MS\_745.9}) were explicitly removed from the feature vocabulary for all patients.
The resulting dataset comprised $N=$12{,}174{,}785 patient vectors. 
Due to the nature of clinical claims data, the feature matrix was highly sparse, containing $N=$489{,}442{,}725 non-zero entries (an overall density of approximately 1.22\%). 
To optimize memory usage and computational efficiency during model training, the data was constructed and stored as a Compressed Sparse Row (CSR) matrix.

%% ----------------------------------------------------------
\subsection{Model Architectures}
%% ----------------------------------------------------------

\subsubsection{BEHRT-LoRA}

BEHRT (BERT for Electronic Health Records)~\cite{TODO_BEHRT} is a
transformer encoder adapted from the original BERT architecture for
medical claims data.
The model was first pre-trained on the full MarketScan patient corpus
using a masked language modelling (MLM) objective, and subsequently
fine-tuned for four-class hip fracture risk prediction.

\paragraph{Encoder architecture.}
The encoder follows the standard BERT design: a stack of transformer
layers, each comprising multi-head self-attention followed by a
position-wise feed-forward network, with residual connections and
layer normalisation at each sub-layer.
The input to the encoder is the summed five-component embedding
described in Section~\ref{sec:features}, passed through layer
normalisation and dropout before the first transformer layer.
The pooled representation of the \texttt{CLS} position after the
final encoder layer is used as the patient-level summary.
The architecture hyperparameters --- hidden dimension, number of
transformer layers, number of attention heads, and feed-forward
intermediate dimension --- were selected by Bayesian hyperparameter
optimisation (Optuna, TPE sampler, 30 trials) on a subset of 100{,}000
training patients over 5 epochs each, searching within the ranges:
hidden dimension $\in \{128, 192, 256, 288, 384\}$; layers~2--8;
attention heads $\in \{4, 6, 8, 12\}$; intermediate dimension
$\in \{256, 512, 768, 1{,}024\}$; dropout~0.05--0.30;
learning rate $10^{-5}$--$5 \times 10^{-4}$.

\paragraph{Pre-training.}
The model was pre-trained on the MarketScan training corpus using the
standard MLM objective: 15\% of tokens per sequence were randomly
selected, of which 80\% were replaced with a \texttt{[MASK]} token,
10\% with a randomly drawn token from the vocabulary, and 10\% left
unchanged.
The model was trained to predict the original token at masked
positions.
Training used a 90/10 train/validation split, with up to 100 epochs
and early stopping (patience 10 evaluations, with 4 evaluations per
epoch).
Mixed-precision training (fp16) and gradient accumulation over 4 steps
were applied, with an effective batch size of $64 \times 4~\text{GPUs}
\times 4~\text{accumulation steps} = 1{,}024$.

\paragraph{Fine-tuning with LoRA.}
Following pre-training, the MLM prediction head was discarded and the
BERT backbone was frozen.
Low-Rank Adaptation (LoRA)~\cite{TODO_LoRA} was injected into the
query and value projection matrices of every attention layer.
LoRA replaces each target weight matrix $W$ with $W + BA$, where $B$
and $A$ are low-rank factors ($A \in \mathbb{R}^{r \times d}$,
$B \in \mathbb{R}^{d \times r}$) initialised so that $BA = 0$ at the
start of fine-tuning.
Only $A$ and $B$ are updated during fine-tuning; the original $W$ is
frozen.
We used rank $r = 16$, scaling factor $\alpha = 32$ (effective scale
$\alpha / r = 2.0$), and LoRA dropout of 0.1, resulting in
approximately 1--2\% of total parameters being trainable.
A linear classification head (dropout followed by a linear projection
to 4 output logits) was appended to the frozen backbone and also
trained.
Fine-tuning used AdamW ($\text{lr} = 2 \times 10^{-4}$, weight decay
0.01) with a linear warmup over 10\% of training steps, gradient
accumulation over 4 steps, and a class-weighted cross-entropy loss
with weights inversely proportional to class frequency.
Training ran for up to 50 epochs with early stopping (patience 10
epochs) on the held-out validation loss.

\subsubsection{XGBoost}

A gradient-boosted decision tree ensemble was trained using
XGBoost~\cite{TODO_XGBoost} on the sparse one-hot binary feature matrix.
XGBoost builds an ensemble of decision trees sequentially, where each
new tree is trained to correct the errors of the trees already in the
ensemble; the final prediction is the weighted sum of all trees'
outputs.

\paragraph{Class imbalance.}
The training set contains approximately 99\% control patients and
$<$1\% fracture cases across the three risk classes.
To prevent the model from ignoring minority classes, the training set
was downsampled so that the control class was capped at 600{,}000
patients (from over 10 million), and sample weights inversely
proportional to class frequency were applied during training.
The validation and test sets were left at their natural class
distributions.

\paragraph{Hyperparameter optimisation.}
Hyperparameters were selected using Optuna Bayesian optimisation
(TPE sampler, 60 trials, 4 trials in parallel across 4 GPUs) with
5-fold stratified cross-validation on the downsampled training set.
The objective minimised mean cross-entropy loss across folds.
The best hyperparameter configuration found was:
maximum tree depth~8; learning rate~0.282; subsampling ratio~0.924;
column subsampling ratio~0.666; minimum child weight~11;
$\gamma = 1.138$; $\alpha = 2 \times 10^{-6}$;
$\lambda = 2.256$; histogram bin count~64.
The final model was trained on the full downsampled training set
using these parameters for 500 boosting rounds with early stopping
on the validation set (patience~50 rounds).

\subsubsection{Feedforward Neural Network}

A feedforward neural network (\texttt{SparseClassifier}) was trained
on the same sparse one-hot feature matrix as XGBoost.
The architecture consists of four fully connected layers with hidden
dimensions 512, 256, and 128, each followed by batch normalisation,
ReLU activation, and dropout.
The final linear layer maps the 128-dimensional hidden representation
to four class logits, from which class probabilities are obtained
via softmax.

\paragraph{Class imbalance.}
The same downsampling strategy as XGBoost was applied: the control
class was capped at 600{,}000 training patients, and a class-weighted
cross-entropy loss was used, with weights inversely proportional to
class frequency in the training set.

\paragraph{Hyperparameter optimisation.}
Hyperparameters were selected using Optuna (TPE sampler, 30 trials)
on a stratified subset of 150{,}000 training patients, with an 80/20
train/validation split per trial and early stopping after 4 epochs
without improvement.
The tuned parameters were learning rate, weight decay, dropout rate,
and batch size.
The best configuration was: learning rate~$2.98 \times 10^{-4}$;
weight decay~$9.59 \times 10^{-4}$; dropout~0.306; batch size~1{,}024.
The final model was trained on the full downsampled training set using
these hyperparameters for up to 60 epochs with early stopping
(patience~10 epochs) based on validation cross-entropy loss.
The AdamW optimiser was used with a \texttt{ReduceLROnPlateau}
learning rate scheduler (factor~0.5, patience~3 epochs).



%% ----------------------------------------------------------
\subsection{External Validation on All of Us}
%% ----------------------------------------------------------

The trained MarketScan models were applied without modification to
the All of Us diabetic cohort ($N=98{,}249$ with EHR data).
Vocabulary and feature alignment procedures are described in
Section~\ref{sec:features}.
Sequence truncation for BEHRT followed the same procedure as
MarketScan: only events prior to the diabetes index date were
included.

Due to the computational cost of transformer inference (approximately
two orders of magnitude greater than scoring sparse feature vectors),
BEHRT was evaluated on a class-stratified random subsample of
$N=50{,}000$ patients.
XGBoost and the neural network were evaluated on the full diabetic
cohort.
Because the BEHRT subsample was stratified by risk class, within-class
proportions are preserved, keeping AUROC, AUPRC, and per-class metrics
directly comparable across all three models.

%% ----------------------------------------------------------
\subsection{Evaluation Metrics}
%% ----------------------------------------------------------

Model performance was assessed using the following metrics, computed
on the held-out test sets for both MarketScan and All of Us.

\paragraph{Overall performance.}
We report macro-averaged precision, recall, F1-score, and
one-vs-rest ROC-AUC.
Macro averaging was preferred over weighted averaging because the
dominant $>$10-year class ($\approx$99\% prevalence) would otherwise
mask meaningful differences in minority-class performance.

\paragraph{Per-class metrics.}
Precision, recall, and F1-score are reported separately for each of
the four time-to-fracture strata.

\paragraph{Diagnostic performance.}
We report the true positive rate (TPR/sensitivity), false positive
rate (FPR), positive predictive value (PPV/precision), negative
predictive value (NPV), and predicted positive rate (PPR) per class.

\paragraph{Calibration.}
Probability calibration was evaluated using the Brier score for each
class and model under three conditions: raw (uncorrected) model
probabilities, Platt scaling (logistic regression recalibration), and
isotonic regression recalibration.
Calibration was assessed on a held-out calibration-test split
comprising 30\% of the All of Us evaluation set.

\paragraph{Confusion matrices.}
Full $4 \times 4$ confusion matrices are reported for each model on
both datasets to characterize patterns of misclassification.

%% ----------------------------------------------------------
\subsection{Statistical Considerations}
%% ----------------------------------------------------------

All performance comparisons between models were conducted on
identical test-set partitions where possible.
For All of Us, BEHRT was evaluated on a stratified 50K subsample
while XGBoost and the neural network were evaluated on the full
cohort; because stratification preserves class proportions, AUROC,
AUPRC, and per-class recall are directly comparable across models,
whereas absolute counts are not.
No formal statistical significance testing was pre-specified for
comparisons between model performance metrics; differences are
reported descriptively.
%% TODO: Add any IRB approval statement, data access agreement
%%       language, or reproducibility/code availability statement
%%       required by the target journal.

%% ============================================================
%%  END OF METHODS SECTION
%% ============================================================


% Bibliography
\bibliographystyle{vancouver}  % or apalike, unsrt, etc.
\bibliography{references}      % references.bib file

\end{document}
